\documentclass[10pt]{article}
\usepackage{icmlko}

\icmlmeta{IP 파운데이션 월드모델}{479}{널이 두 종류였다}{어제 세운 자가 잡음 바닥 아래라는 지적을 확인하러 갔더니 --- 지적은 옳았고, 근거는 죽은 숫자였고, 답은 문턱을 올리는 것이 아니었다}

\begin{document}

\twocolumn[
\icmlheader

\begin{abstract}
직전 사이클이 채택 문턱을 $0.00353$ 으로 세웠다. 프레시 티처가 하루 만에
반박했다 --- \emph{``이 실험실이 잰 잡음 바닥(순수 난수 3열 $=0.0043$)보다
낮다.''} 그 지적을 확인하러 갔다.

\textbf{지적은 옳았다. 다만 근거로 인용된 $0.0043$ 은 죽은 숫자였다} ---
노트 742 가 이미 $3.7\sigma$ 이상치로 격하했고 산 값은 $0.0153$ 이다.
그리고 그 정정은 저장소의 은퇴 숫자 검출기에 \textbf{등록돼 있었다}.
오늘 검출기를 돌리면 \emph{티처가 적은 인계 카드 그 줄이 걸린다}.

12도메인 $\cdot$ 유보 3,775 에서 \textbf{뽑기 12회}로 다시 쟀다. 신호가
$0$ 인 연속 난수 열 \textbf{하나}를 넣으면 판이 \textbf{$+0.00558$} 내려가고
(12뽑기 \textbf{전부} 같은 부호), 바닥 $\mu{+}2\sigma$ 는 \textbf{$0.01055$}
--- 채택 문턱의 \textbf{2.99배}다. 3열이면 \textbf{$0.02126$(6.02배)}.

그래서 답은 \emph{``문턱을 올려라''} 가 아니었다. \textbf{널이 두
종류다} --- 열 수가 그대로인 개입과 열이 느는 개입은 참 $\Delta$ 자체가
다르다. 앞의 널은 $0$ 이고 뒤의 널은 $0$ 이 아니다. 열이 느는 개입에
$0.00353$ 을 대면 \textbf{3배 헐겁게 읽는다}.
\end{abstract}
]

\section{하루 만의 반박}

직전 사이클은 53노트 묵은 미결을 이행해 채택 문턱을 $0.0045$ 에서
$0.00353$ 으로 바꿨다. 그날 밤 프레시 티처(이 사이클을 돌지 않은 세션)가
치명 넷을 냈고 그중 첫째가 이것이었다:

\begin{quote}
\emph{``새 문턱이 이 실험실이 잰 잡음 바닥 아래로 내려갔다 --- 노트 146 이
순수 난수 연속 3열을 $0.0043$ 으로 실측했는데 새 문턱은 그 $0.82$배다.
신호 $0$ 인 열 셋이 채택 문턱을 넘는다.''}
\end{quote}

\textbf{채택 문턱이 자기 실험실의 잡음 바닥보다 낮으면 그것은 채택 문턱이
아니다.} 논증은 옳다. 그래서 이 사이클은 그 바닥을 12도메인에서 다시 쟀다.

\claimx{바닥을 다시 재면 문턱과의 관계가 정해진다.}{바닥이 문턱 아래로
나오면 티처의 지적이 철회되고, 위로 나오면 문턱이 무효다.}

\section{인용된 수가 죽어 있었다}

측정 전 사전등록 단계에서 인용값을 먼저 세었다(규율: \emph{예측을 적기 전에
그 값이 이미 파일에 있는지 확인한다}). 노트 146 의 $0.0043$ 은
\textbf{노트 742 가 $3.7\sigma$ 이상치로 격하}했고 산 값은 $0.0153$ 이다.

더 나쁜 것은 그 정정이 \textbf{이미 자동 검출기에 등록돼 있었다}는 점이다.
`ingest/audit.py` 의 은퇴 숫자 목록에 그 항목이 있고, 오늘
`dead\_numbers()` 를 돌리면 \textbf{티처가 방금 적은 인계 카드 그 줄이
걸린다.} 검출기가 있고, 항목이 등록돼 있고, 아무도 안 돌렸다.

\claimx{인용은 세지 않으면 죽은 채로 옮겨 다닌다 --- 검출기가 있어도.}{그
검출기가 그 줄을 안 잡으면 검출기 쪽이 고장이다.}

\section{어떻게 쟀나}

\textbf{배선을 먼저 닫았다(여섯).} 자료 sha256 `d7e2df86'(890$\cdot$891 과
동일) $\cdot$ 유보 3,775 $\cdot$ \textbf{`axis\_order' $36 \to 39$}(난수 열이
\emph{실제로} 설계행렬에 들어갔는지) $\cdot$ 연속 열 고유값이 행수의
99.9\% 이상 $\cdot$ 마스크 전부 1(중립화 없음) $\cdot$ 포개기 확인.

그리고 \textbf{항등 통제}를 넣었다 --- 같은 팔을 같은 씨앗으로 재구축하면
$\Delta$ 가 \emph{부동소수 정확히} $0.0$ 인가. 둘 다 통과했다. 직전 사이클이
빠뜨린 한 줄이고, 노트 887 이 무효화된 바로 그 자리다.

\textbf{뽑기 12회.} 노트 146 은 난수 열을 \textbf{한 번} 뽑았다. 바닥은
분포의 분위수이지 한 번 뽑은 값이 아니므로(규약 20 확장), 뽑기 12 $\times$
씨앗 6 으로 다시 쟀다. 적합 302회 $\cdot$ 1,997초.

\section{바닥은 문턱의 세 배였다}

\begin{figure*}[t]
\centering
\includegraphics[width=\textwidth]{figs/fig1.pdf}
\caption{왼쪽: 신호가 $0$ 인 열을 넣으면 판이 내린다 --- 바닥(막대)과 뽑기
평균$\pm\sigma$(점), 붉은 선이 채택 문턱. 가운데: 거칠기 사다리는 단조다.
오른쪽: 얇은 도메인은 \emph{평균}이 아니라 \emph{분산}이 크다.}
\label{fig:1}
\end{figure*}

\begin{center}\small
\begin{tabular}{lrrrr}
\toprule
열 수 & 뽑기 $\mu$ & 뽑기 $\sigma$ & 바닥 $\mu{+}2\sigma$ & 문턱 대비 \\
\midrule
1열 & $+0.00558$ & 0.00248 & \textbf{0.01055} & \textbf{2.99$\times$} \\
2열 & $+0.00977$ & 0.00379 & \textbf{0.01735} & \textbf{4.92$\times$} \\
3열 & $+0.01335$ & 0.00396 & \textbf{0.02126} & \textbf{6.02$\times$} \\
\bottomrule
\end{tabular}
\end{center}

\textbf{12뽑기가 전부 같은 부호다} --- 열 하나에서도 열 셋에서도 판은
\emph{언제나} 내려간다. $\mu$ 구간은 `cluster\_boot' BCa 로 냈다(규약 47).
그리고 거칠기 사다리가 단조다: $k{=}2\ 0.00241 \to k{=}4\ 0.00391 \to
k{=}10\ 0.00710 \to k{=}100\ 0.01157 \to$ 연속 $0.01286$($k{\approx}100$
포화). 노트 146 의 정성적 그림은 살아남았고 눈금만 바뀌었다.

\claimx{신호가 하나도 없는 열도 판을 내린다 --- 그리고 그 크기는 채택 문턱의
세 배다.}{12뽑기 중 셋 이상이 반대 부호면 이것은 잡음이지 비용이 아니다
--- 실측은 $12/12$ 였다.}

\section{그래서 답은 문턱이 아니었다}

문턱을 $0.02126$ 으로 올리면 될까. 안 된다. 그러면 \emph{열 수를 안 바꾸는}
개입(손잡이 조정 $\cdot$ 치환 $\cdot$ 자료 고침)까지 6배 엄한 자로 재게 된다.

\textbf{두 개입은 널 자체가 다르다.}

\begin{center}\small
\begin{tabular}{lll}
\toprule
 & 열 고정 개입 & 열 증가 개입 \\
\midrule
참 $\Delta$ & $0$ & \textbf{$0$ 이 아니다} \\
널의 정체 & 두 팔이 같을 때 & 신호 $0$ 인 같은 꼴 \\
자 & 0.00353 & 0.01055 (1열) \\
 & & 0.02126 (3열) \\
\bottomrule
\end{tabular}
\end{center}

직전 사이클의 $0.00353$ 은 \textbf{틀리지 않았다 --- 적용 범위가 안 적혀
있었을 뿐이다.} 열이 느는 개입에 그 자를 대면 3배 헐겁게 읽는다.

이 눈으로 원장을 다시 훑으니 $0.0045 \to 0.00353$ 으로 편이 바뀌는 항목이
\textbf{8건}, $\to$ 3열 바닥으로는 \textbf{77건}이다. 직전 사이클의
\emph{``뒤집힘 0건''} 은 자기 사이클만 본 것이었다.

\claimx{자는 하나가 아니라 개입의 꼴에 달렸다.}{두 널의 $2\sigma$ 가
겹치면 하나로 써도 된다 --- 실측은 $2.99$--$6.02$배로 안 겹쳤다.}

\section{빗나간 예측 하나가 제일 값졌다}

사전등록 여덟 중 일곱이 맞았다. 틀린 하나는 P8 --- \emph{``$|$흔들림$|$ 과
유보 행 수는 음의 상관''}(노트 146 의 \emph{``얇은 도메인이 크게 흔들린다''}).
실측 스피어만 $+0.126$($p{=}0.70$) 으로 \textbf{빗나갔다}.

사후 탐색이 자리를 찾아 줬다. 음의 상관은 \textbf{평균이 아니라 뽑기 사이
SD} 에 붙는다 --- $\rho = -0.692$($p{=}0.013$). 시장팝업(유보 126)은 평균
흔들림 $-0.0643$ 인데 \textbf{뽑기 SD 가 $0.0908$} 로 평균보다 크다.

\textbf{얇음은 판을 한쪽으로 밀지 않는다. 판을 예측 불가능하게 만든다.}
같은 도메인이 뽑기에 따라 크게 오르기도 크게 내리기도 한다. 이것은 사후
발견이므로 \textbf{다음 사이클의 사전등록 대상}이다.

\section{남는 것}

이 사이클도 채택을 안 냈다 --- 사전등록이 \emph{``첫 양수를 그대로 채택하지
않는다''}(노트 133)를 명시했고 그대로 지켰다. 새 바닥은 \textbf{제안}이고
정본 승격은 다음 사이클의 확인 측정 뒤다.

그리고 아직 배선이 없다. \textbf{널의 종류가 판정 규칙에 안 들어가 있다} ---
지금은 인계 카드와 README 에만 있다. 코드가 두 널을 갈라 쓰기 전까지 이
사이클의 결론은 문서에만 산다. 직전 사이클이 정확히 그 상태로 하루를
보냈고, 그래서 티처가 죽은 숫자를 인용할 수 있었다.

\end{document}
